% ==============================================================================
% SECCIÓN 12: PLAN DE MEDICIÓN DEL DESEMPEÑO, RECOMPENSAS Y CASTIGOS
% ==============================================================================
\section{Plan de Medición del Desempeño, Recompensas y Castigos}

\subsection{Registro y Control del Esfuerzo Individual}
Cada miembro cumplimentará semanalmente su hoja personal de registro de dedicación horaria en la plantilla oficial (\texttt{docs/plantillas/plantilla\_registro\_horas.md}). El balance consolidado será auditado por el Coordinador y el Gestor de Calidad en cada iteración para garantizar un reparto equilibrado de la carga.

\subsection{Sistema de Recompensas y Reconocimiento}
\begin{itemize}[leftmargin=*]
  \item \textbf{Reconocimiento Público Interno:} En el acta de retrospectiva de cada iteración se hará mención explícita al miembro o rol destacado por su proactividad, aportes de calidad o apoyo extraordinario a sus compañeros.
  \item \textbf{Prioridad de Elección Técnica:} Aquellos integrantes con desempeño excelente y cumplimiento anticipado de entregables tendrán preferencia al seleccionar sus componentes de trabajo en el sprint sucesivo.
  \item \textbf{Objetivo Común de Calificación Máxima:} Se busca promover la excelencia compartida para que todos los integrantes obtengan la máxima calificación académica (10) en la evaluación de prácticas.
\end{itemize}

\subsection{Régimen Sancionador Interno (Castigos y Penalizaciones)}
Para asegurar la responsabilidad individual y salvaguardar los resultados del grupo frente a actitudes negligentes:
\begin{enumerate}[leftmargin=*]
  \item \textbf{Inasistencia a Reuniones:} La falta injustificada a una reunión conllevará apercibimiento interno. La acumulación de \textbf{dos faltas injustificadas} a reuniones oficiales acarreará amonestación formal en acta y se traducirá en penalización en la nota individual remitida a los profesores responsables.
  \item \textbf{Incumplimiento de Entregables o Plazos:} El retraso o no entrega de tareas comprometidas sin previo aviso justificado de al menos \textbf{48 horas} motivará la reasignación de urgencia de la tarea y figurará con calificación desfavorable en el informe individual del implicado.
  \item \textbf{Reincidencia en Falta de Calidad:} Subir código defectuoso que impida el trabajo del resto obligará a su corrección inmediata fuera del cómputo computable de horas.
\end{enumerate}

\vfill
\begin{center}
  \small\sffamily\color{textGray}
  Documento consensuado y aprobado por el Equipo CLENCH Software Development.\\
  Grado en Ingeniería Informática --- Escuela Técnica Superior de Ingenierías Informática y de Telecomunicación.\\
  Universidad de Granada, Curso Académico 2026--2027.
\end{center}
